\begin{abstract}
This paper empirically investigates whether exogenous alternative data---specifically, natural language processing (NLP) derived news sentiment and market microstructure signals---contains incremental predictive power for forecasting realized volatility beyond the classical Heterogeneous Autoregressive (HAR) framework. We conduct a cross-asset evaluation using high-frequency Bitcoin (BTC), the S\&P 500 (SPX), and the NIFTY 50 index. We compare a linear sentiment-augmented model (HAR-S) against a non-linear deep learning hybrid (Neural-HAR) that employs a Gated Residual Network and an FT-Transformer to predict the residual error of an ordinary least squares HAR prior. We find that linear models struggle to extract out-of-sample signal from high-dimensional sentiment data. In contrast, the Neural-HAR model yields statistically significant improvements in forecasting accuracy, primarily under the QLIKE loss function, and acts as the sole surviving model in the Model Confidence Set for SPX. However, we note an important limitation: the Neural-HAR evaluation is conducted on a shorter out-of-sample window ($\approx$ 1 trading year) compared to the multi-year baseline HAR evaluation. In an economic application using a volatility-targeting strategy, the superior forecasts of the Neural-HAR model do not improve the Sharpe ratio, but they systematically reduce realized portfolio volatility and maximum drawdowns, offering enhanced risk control. Finally, we find no robust evidence that sentiment's marginal predictive content is strictly regime-dependent.
\end{abstract}
